\documentclass[11pt]{article}
\usepackage[margin=0.9in]{geometry}
\usepackage[T1]{fontenc}
\usepackage{amsmath,amssymb,amsthm,booktabs,array}
\usepackage[colorlinks=true,urlcolor=blue,linkcolor=blue]{hyperref}
\newtheorem{theorem}{Theorem}
\newtheorem{lemma}{Lemma}
\newcommand{\F}{\mathbb F_p}
\newcommand{\one}{\mathbf1}
\title{\textbf{PrismGuard-32}\\Masked Key Capsules with Bounded Recovery\\and an Executable Probing Laboratory}
\author{Mathematics for Systems Security --- Cryptographer Track}
\date{October 25, 2026}
\begin{document}
\maketitle
\begin{abstract}
PrismGuard-32 is an experimental master-key capsule and protected-computation
case study. Its contribution is a specific integration of randomized
two-share arithmetic, a masked acceptance predicate, an eight-attempt
receiver state machine, and reproducible attack experiments. A custom
finite-field capsule supplies an inspectable one-use confidentiality
foundation; no novelty or additional hardness is attributed to its public
matrix and power maps. We give correctness and ideal public-adversary
substitution arguments, analyze a single-field-wire arithmetic model, and
exhaustively check the shared multiplication schedule over a small field
in Rust and independently in Python. Deliberately unsafe variants expose
the boundary of protection. The results do not establish a new multi-use
encryption primitive, whole-program formal verification, compiler-level
noninterference, or physical side-channel resistance.
\end{abstract}

\section{Research question and contribution}
Can a key-capsule application avoid unsharing private nonlinear computations,
verify tags without exposing raw differences, and tolerate several hostile
packets before legitimate recovery? PrismGuard-32 addresses this narrow
question with source, an in-memory application, proofs under explicit
assumptions, and negative controls. The preserved PrismLock-32 baseline
used raw private arithmetic and allowed only one receiver attempt.
The new implementation shares that arithmetic and changes the recovery
contract, with a corresponding change to the probability bound.

Finite fields, full-entropy masking, power permutations, ISW-style
masking~\cite{isw}, and polynomial authentication~\cite{cw} are established
techniques. The exact application, composition, transcript, release policy,
state contract, and evidence harness were developed for this entry.
We do not claim to invent the mathematical foundations or prove worldwide
novelty. The organizer must interpret the track's restriction on existing
plaintext-encryption algorithms; direct implementation alone does not
resolve that interpretation.

\section{Parameters and wire protocol}
Let
\[
p=2^{61}-1,\quad
A=I+\one\one^T\in\F^{8\times8},\quad
A^{-1}=I-9^{-1}\one\one^T.
\]
Here $9^{-1}=2049638230412172401\bmod p$, $\det A=9$, and
\[
d=1763291712928118903,\qquad 17d\equiv1\pmod{p-1}.
\]
The coordinate power $x\mapsto x^{17}$ is a permutation, including at zero.
The independent reference checks the prime with the Lucas--Lehmer criterion
for prime exponent 61, obtaining residue zero after 59 iterations.

Encode a 32-byte master key as eight little-endian 32-bit words
$m\in\{0,\ldots,2^{32}-1\}^8\subset\F^8$.
Fresh logical material consists of
\[
r\gets\F^8,\quad
\alpha_0,\alpha_1\gets\F^\ast,\quad
\beta_0,\beta_1\gets\F.
\]
A public 16-byte identifier and public context $a$ of at most 256 bytes
identify the capsule and application. Text contexts use UTF-8.
The capsule's logical encryption is
\begin{equation}
 c=A\bigl((Am+r)^{[17]}\bigr).
 \label{eq:capsule}
\end{equation}
The power is coordinatewise. Full-entropy masking supplies secrecy;
the public bijections add no entropy or cryptographic hardness.

Let $B(b)$ split bytes into consecutive chunks of at most seven bytes,
each interpreted little-endian as a field value. Define the injective
transcript
\[
T=(\mathtt{0x504755415244},\,2,\,B(\mathrm{id}),\,|a|,\,B(a),\,
c_0,\ldots,c_7).
\]
There are at most $1+1+3+1+37+8=51$ coefficients. The context byte length
prevents ambiguity in its last zero-padded chunk. For $T=(t_1,\ldots,t_N)$,
\begin{equation}
H_\alpha(T)=\alpha^{N+1}+\sum_{j=1}^{N}t_j\alpha^{N+1-j},\quad
\tau_\ell=H_{\alpha_\ell}(T)+\beta_\ell \pmod p.
\label{eq:tag}
\end{equation}
The Horner implementation begins at 1 and finishes with an additional
multiplication by $\alpha$. Omitting the final multiplication allows a
known additive modification to the final ciphertext coefficient and tag.

\begin{center}
\begin{tabular}{rll}
\toprule
Offset & Bytes & Meaning\\
\midrule
0 & 8 & \texttt{PGUARD}, zero byte, version byte 2\\
8 & 16 & Public identifier\\
24 & 64 & Eight canonical ciphertext words\\
88 & 16 & Two canonical tag words\\
\bottomrule
\end{tabular}
\end{center}
All field words are unsigned 64-bit little-endian values below $p$.
Context is supplied separately and authenticated. The format is exactly
104 bytes. Noncanonical values, wrong magic/version, incorrect length,
and overlong contexts are rejected.

\section{Protected arithmetic}
A private field value $x$ is represented as $(x_L,x_R)$ with
$x=x_L+x_R\bmod p$. Trusted ingress and provisioning perform the initial
encoding. The protected phase never explicitly reconstructs a private
logical value except at a documented public-output or authorized-release
boundary. Both shares are in the same process; this is not physical
domain separation.

\subsection{Refresh and linear operations}
For independent uniform $\delta$,
\[
\operatorname{Refresh}(x_L,x_R)=
(x_L+\delta,\ x_R-\delta).
\]
Addition and subtraction refresh one operand first. This is necessary:
blindly adding correlated representations can cancel masks.
Public addition changes one component; public scaling scales both.

\subsection{Multiplication, including correlated inputs}
Refresh both operands independently, even for squaring. For logical values
$a,b$, let the refreshed first components be independent uniform $U,V$,
with second components $a-U,b-V$. Take independent uniform $\rho$ and compute
\begin{align*}
w_{00}&=UV,&w_{01}&=U(b-V),\\
w_{10}&=(a-U)V,&w_{11}&=(a-U)(b-V),\\
z_L&=w_{00}+\rho,&q_0&=w_{01}-\rho,\\
q_1&=q_0+w_{10},&z_R&=q_1+w_{11}.
\end{align*}
The order matters: no unmasked accumulation of the three cross terms occurs.
The implementation uses the same generic source schedule for the production
field and the small-field audit.

\begin{lemma}[Arithmetic correctness]
The output represents $ab$.
\end{lemma}
\begin{proof}
Summing its components cancels $\rho$ and includes exactly the four products
in $(U+(a-U))(V+(b-V))=ab$.
\end{proof}

\begin{lemma}[Single-wire gadget distributions]
For fixed logical $a,b$ and fresh independent input refreshes and $\rho$,
every observed private field wire in the stated multiplier schedule has a
distribution independent of $a,b$, assuming input share marginals have that
property.
\end{lemma}
\begin{proof}
The refreshed $U,V$ are independent uniform field elements even when the
logical inputs coincide. Each $w_{ij}$ is distributed as a product of two
independent uniform field elements: its mass at zero is $(2p-1)/p^2$ and
at any nonzero value is $(p-1)/p^2$, independent of $a,b$.
The output and accumulator values masked by $\rho$ are uniform.
Individual old input shares have the assumed independent marginals;
refresh deltas are independent randomness; individual refreshed components
are uniform. Thus each listed wire has an input-independent distribution.
This is a statement about the specified arithmetic schedule.
\end{proof}

The ideal circuit argument follows the same conditional-distribution
reasoning through refreshed binary gates: logical values are unaffected
by blinding coins, and each private share/temporary has an input-independent
marginal. Public constants and deliberately released outputs are treated
as public. We do not claim that the small-field audit constitutes a formal
whole-program proof, nor infer robust or higher-order composability from
single-wire marginals. Physical defaults and composability require stronger
models~\cite{robust}.

\subsection{A deliberately unsafe alternative}
If a multiplier forms the unmasked sum of its three cross terms, that value
is $ab-UV$. Its distribution is a shifted product distribution and depends
on $ab$. Over $\mathbb F_5$, the exact experiment finds total-variation
distance $0.20$ between the selected input cases. Both output shares
together reconstruct $ab$, demonstrating the excluded two-probe boundary.

\section{Masked verification and controlled release}
Compute both expected tags entirely in shares. Subtract the public candidate
tag in the shared representation, without releasing either difference.
For a difference $\Delta$ compute
\[
b=\Delta^{p-1}=
\begin{cases}0,&\Delta=0,\\1,&\Delta\ne0.\end{cases}
\]
Exponentiation uses shared multiplication with a fixed public exponent.
Combine both results in shares:
\[
b_{\rm reject}=b_0+b_1-b_0b_1.
\]
Only this one-bit predicate is released. Public parser checks may fail
earlier. After successful authentication, compute
\[
 m=A^{-1}\left((A^{-1}c)^{[d]}-r\right)
\]
in shares, then cross the authorized plaintext boundary. Reject if a
recovered field word does not fit in 32 bits, and release no partial
plaintext to the caller. The returned secret has an owned wipe path.

Explicit release sites in the source are limited to public ciphertext,
public authentication tags, the authentication predicate, and authenticated
plaintext. Timing of success and rejection paths need not match:
the outcome is intentionally public. A generic rejection error avoids
exposing a finer tag-failure oracle.

\section{Secrecy and bounded public substitution}
\begin{theorem}[Ideal one-use public confidentiality]
For ideal uniform independent logical material, the public capsule
distribution is independent of $m$ for fixed public identifier/context.
\end{theorem}
\begin{proof}
For each fixed $m$, translation by $Am$ maps uniform $r$ to uniform
$Am+r$. Both subsequent maps in~\eqref{eq:capsule} are bijections of
$\F^8$. Independently uniform $\beta_0,\beta_1$ make both tags uniform.
Thus every message induces the same public distribution.
\end{proof}
The implementation uses OS-CSPRNG output, so this ideal statement does not
make the practical randomness information-theoretically random.

\begin{lemma}[One public substitution]
After one genuine capsule is observed, any one altered well-formed
transcript is accepted with probability at most
$\epsilon=(52/(p-1))^2$ in the no-additional-leakage public-adversary model.
\end{lemma}
\begin{proof}
Each fresh $\beta_\ell$ leaves $\alpha_\ell$ uniform and hidden after the
valid tag is observed. Different injective encodings yield a nonzero
polynomial difference of degree at most 52 with zero constant term.
Subtracting any chosen tag difference cannot make that polynomial zero
identically. It has at most 52 roots in $\F^\ast$.
The two evaluation points are independent, giving the product bound.
An unchanged transcript with a different tag is rejected deterministically.
\end{proof}

\begin{theorem}[At most eight public attempts]
For the receiver state machine, the probability of any altered-transcript
acceptance in $Q\le8$ public verification attempts is at most
\[
Q\epsilon\le8\left(\frac{52}{p-1}\right)^2
\approx4.069\cdot10^{-33}\approx2^{-107.6}.
\]
\end{theorem}
\begin{proof}
Fix the adversary's external randomness and observed public capsule.
Before the first acceptance, its all-rejection path determines at most
$Q$ candidate forgeries. Each candidate's acceptance event obeys the
one-substitution bound. A union bound applies to these events, with no
independence assumption between attempts. In particular, we do not
incorrectly assume that the hidden points remain freshly uniform after
conditioning on each rejection.
\end{proof}
An unchanged genuine capsule is not a forgery. Numerical authentication
bounds above exclude additional key-dependent leakage; the single-wire
circuit analysis is stated separately. No unrestricted joint leakage/CCA
theorem is claimed.

\section{Receiver state and availability}
Provisioning creates a non-\texttt{Clone}, non-\texttt{Copy} sender handle
and receiver object with independent share representations of identical
logical material. Sealing consumes the sender handle. The receiver permits
eight attempts:
\[
\mathrm{Ready}(8)\longrightarrow\cdots\longrightarrow
\begin{cases}
\mathrm{Opened},&\text{valid recovery},\\
\mathrm{Locked},&\text{eighth rejection}.
\end{cases}
\]
Success consumes logical material. The eighth rejection discards it.
Malformed packets consume the same attempt budget. Insufficient blinding
entropy returns a resource error before an attempt starts.
Seven bad packets can therefore precede a valid eighth recovery.
Eight bad packets still deny service. In-memory state does not prevent
snapshots, unsafe copies, or rollback; no durable ledger is implemented.

Never encrypt a second message under the same logical pad. Define
$Z(c)=(A^{-1}c)^{[d]}$. Under reuse,
\[
A^{-1}(Z(c')-Z(c))=m'-m.
\]
A known plaintext reveals the second plaintext. The Python reference
contains a successful synthetic attack exhibiting this excluded use.

\section{Implementation and measurements}
Production code uses no third-party crates or existing plaintext-cipher
implementation. Field reduction uses Mersenne folding and conditional
arithmetic rather than division. Binary masking gates are kept as explicit
function boundaries to make the target easier to inspect; this is not a
compiler-preservation proof. Private-dependent table indices are absent.
Owned buffers have volatile wipe paths. Moves, spills, registers, caller
copies, paging and crash dumps remain outside a complete-erasure guarantee.

Logical material uses 96 bytes; two-share storage takes 192 bytes per
endpoint plus the public identifier. A masked master key takes 128 bytes.
The implementation provisions a 64 KiB blinding pool, of which the maximum
observed seal/open consumption at context size 256 is 482/3317 field words.

\begin{center}
\begin{tabular}{lrr}
\toprule
Recorded operation & Samples & Median\\
\midrule
Seal, pool provisioning excluded & 1000 & $1.7\,\mu$s\\
Valid open, pool provisioning excluded & 1000 & $18.8\,\mu$s\\
Full workflow including OS RNG & 200 & $88.5\,\mu$s\\
\bottomrule
\end{tabular}
\end{center}
The complete raw summary and compiler/source hashes accompany the entry.
Results are from one Windows x64 build, not performance promises.

\section{Verification record}
\begin{itemize}
\item Fourteen Rust test groups; all 832 individual capsule-bit modifications
rejected; all byte values roundtripped; context/parser/state boundaries checked.
\item 512 full capsule vectors matched by independent Python big-integer code;
60,000 field-operation results checked.
\item 2000 reproducible multi-bit mutations rejected, with valid recovery after
each rejection. This is not a coverage-guided fuzzing claim.
\item The exact generic multiplier schedule runs 78,125 times over $\mathbb F_5$:
25 logical pairs, 3125 coin assignments each, 19 observed wires. All 475
per-input wire-distribution checks agree. Python independently reproduces them.
\item Unsafe cross-term accumulation is distinguishable; two shares recover the
logical output; a reused-pad known-plaintext attack succeeds as expected.
\item Compile-fail checks reject sender reuse and receiver cloning.
\item Three timing runs cover seal, valid open, rejected open, and an intentional
large-leak negative control, with 10,000 samples per class per case.
Raw and pooled-percentile-cropped Welch statistics are retained.
\item This run has no exploratory protected-case statistic above 4.5. Negative
controls are detected. This is not a constant-time certificate or a calibrated
whole-experiment significance claim; testing limitations apply~\cite{dudect}.
\item Saved optimized library assembly contains no integer division instruction.
The inspected refresh/add/subtract/multiply bodies have no conditional jump;
entropy access has public bounds checks. No formal binary dataflow proof is made.
\end{itemize}

\section{Trust boundaries, disclosure, and remaining work}
Trusted ingress/provisioning handles raw values. Authorized release returns
the recovered master key. The arithmetic result does not cover two probes,
continual leakage collected across executions, transition/glitch leakage,
power/EM, active faults, speculative execution, compromised operating systems,
or rollback. Physical leakage and compiler-preservation proofs require
additional work; single-wire reasoning must not be relabeled as hardware
resistance~\cite{intel,robust}.

All testing uses synthetic data in a newly created folder. No existing user
files are encrypted, edited or deleted. The Windows demo performs no file,
network, registry, system-setting, or subprocess actions; its external service
is operating-system randomness. The verification harness invokes the installed
compiler and writes its own build/evidence files. Baseline hashes are compared
before and after verification.

AI assistance was used for design, implementation, testing and writing.
No independent human cryptographic audit or physical test campaign has occurred.
Organizer acceptance and winning cannot be guaranteed.

\begin{thebibliography}{9}
\bibitem{isw}Y. Ishai, A. Sahai, D. Wagner,
\emph{Private Circuits: Securing Hardware against Probing Attacks}, CRYPTO 2003.
\url{https://iacr.org/archive/crypto2003/27290462/27290462.pdf}
\bibitem{robust}S. Faust et al.,
\emph{Composable Masking Schemes in the Presence of Physical Defaults and the
Robust Probing Model}, TCHES 2018.
\url{https://eprint.iacr.org/2017/711}
\bibitem{cw}M. N. Wegman, J. L. Carter,
\emph{New classes and applications of hash functions}, FOCS 1979.
\url{https://research.ibm.com/publications/new-classes-and-applications-of-hash-functions}
\bibitem{dudect}O. Reparaz, J. Balasch, I. Verbauwhede,
\emph{Dude, is my code constant time?}, DATE 2017; implementation and limitations.
\url{https://github.com/oreparaz/dudect}
\bibitem{intel}Intel, \emph{Guidelines for Mitigating Timing Side Channels Against
Cryptographic Implementations}.
\url{https://www.intel.com/content/www/us/en/developer/articles/technical/software-security-guidance/secure-coding/mitigate-timing-side-channel-crypto-implementation.html}
\end{thebibliography}
\end{document}
